\documentclass[10pt,a4paper]{amsart}
\usepackage{amsmath,amssymb,amsthm,graphicx}
% Retained minimal preamble for the Stacks Project chapter varieties.tex.
% Verbatim excerpt of the chapter preamble: Xy-pic setup (native preamble lines 14-18), the native theorem
% declarations with the mathematical macros (native preamble lines 154-217), and the project own no-op
% definitions of the reference/slogan/history display environments (native wrapper.tex lines 1-3) that the
% counted statement uses for its slogan. Copied mechanically; nothing added, removed or redefined.
% For commutative diagrams we use Xy-pic
\usepackage[all]{xy}

% We use 2cell for 2-commutative diagrams.
\xyoption{2cell}
% Theorem environments.
%
\theoremstyle{plain}
\newtheorem{theorem}[subsection]{Theorem}
\newtheorem{proposition}[subsection]{Proposition}
\newtheorem{lemma}[subsection]{Lemma}

\theoremstyle{definition}
\newtheorem{definition}[subsection]{Definition}
\newtheorem{example}[subsection]{Example}
\newtheorem{exercise}[subsection]{Exercise}
\newtheorem{situation}[subsection]{Situation}

\theoremstyle{remark}
\newtheorem{remark}[subsection]{Remark}
\newtheorem{remarks}[subsection]{Remarks}

\numberwithin{equation}{subsection}

% Macros
%
\def\lim{\mathop{\mathrm{lim}}\nolimits}
\def\colim{\mathop{\mathrm{colim}}\nolimits}
\def\Spec{\mathop{\mathrm{Spec}}}
\def\Hom{\mathop{\mathrm{Hom}}\nolimits}
\def\Ext{\mathop{\mathrm{Ext}}\nolimits}
\def\SheafHom{\mathop{\mathcal{H}\!\mathit{om}}\nolimits}
\def\SheafExt{\mathop{\mathcal{E}\!\mathit{xt}}\nolimits}
\def\Sch{\mathit{Sch}}
\def\Mor{\mathop{\mathrm{Mor}}\nolimits}
\def\Ob{\mathop{\mathrm{Ob}}\nolimits}
\def\Sh{\mathop{\mathit{Sh}}\nolimits}
\def\NL{\mathop{N\!L}\nolimits}
\def\CH{\mathop{\mathrm{CH}}\nolimits}
\def\proetale{{pro\text{-}\acute{e}tale}}
\def\etale{{\acute{e}tale}}
\def\QCoh{\mathit{QCoh}}
\def\Ker{\mathop{\mathrm{Ker}}}
\def\Im{\mathop{\mathrm{Im}}}
\def\Coker{\mathop{\mathrm{Coker}}}
\def\Coim{\mathop{\mathrm{Coim}}}

% Boxtimes
%
\DeclareMathSymbol{\boxtimes}{\mathbin}{AMSa}{"02}

%
% Macros for moduli stacks/spaces
%
\def\QCohstack{\mathcal{QC}\!\mathit{oh}}
\def\Cohstack{\mathcal{C}\!\mathit{oh}}
\def\Spacesstack{\mathcal{S}\!\mathit{paces}}
\def\Quotfunctor{\mathrm{Quot}}
\def\Hilbfunctor{\mathrm{Hilb}}
\def\Curvesstack{\mathcal{C}\!\mathit{urves}}
\def\Polarizedstack{\mathcal{P}\!\mathit{olarized}}
\def\Complexesstack{\mathcal{C}\!\mathit{omplexes}}
% \Pic is the operator that assigns to X its picard group, usage \Pic(X)
% \Picardstack_{X/B} denotes the Picard stack of X over B
% \Picardfunctor_{X/B} denotes the Picard functor of X over B
\def\Pic{\mathop{\mathrm{Pic}}\nolimits}
\def\Picardstack{\mathcal{P}\!\mathit{ic}}
\def\Picardfunctor{\mathrm{Pic}}
\def\Deformationcategory{\mathcal{D}\!\mathit{ef}}
\newenvironment{reference}{}{}
\newenvironment{slogan}{}{}
\newenvironment{history}{}{}
\usepackage{cleveref}
\crefname{originaldefinition}{original definition}{original definitions}
\Crefname{originaldefinition}{Original definition}{Original definitions}
\makeatletter
\newlabel{morphisms-definition-normalization-X-in-Y}{{035H}{source}{}{sourcepointer-1}{}}
\newlabel{morphisms-definition-normalization-X-in-Y@cref}{{[originaldefinition][0][]035H}{[1][0][]source}{}{}{}}
\makeatother
\crefname{originallemma}{original lemma}{original lemmas}
\Crefname{originallemma}{Original lemma}{Original lemmas}
\makeatletter
\newlabel{morphisms-lemma-finite-type-noetherian}{{01T6}{source}{}{sourcepointer-2}{}}
\newlabel{morphisms-lemma-finite-type-noetherian@cref}{{[originallemma][0][]01T6}{[1][0][]source}{}{}{}}
\makeatother
\crefname{originallemma}{original lemma}{original lemmas}
\Crefname{originallemma}{Original lemma}{Original lemmas}
\makeatletter
\newlabel{morphisms-lemma-nagata-normalization-finite-general}{{0AVK}{source}{}{sourcepointer-3}{}}
\newlabel{morphisms-lemma-nagata-normalization-finite-general@cref}{{[originallemma][0][]0AVK}{[1][0][]source}{}{}{}}
\makeatother
\crefname{originallemma}{original lemma}{original lemmas}
\Crefname{originallemma}{Original lemma}{Original lemmas}
\makeatletter
\newlabel{morphisms-lemma-ubiquity-nagata}{{035B}{source}{}{sourcepointer-4}{}}
\newlabel{morphisms-lemma-ubiquity-nagata@cref}{{[originallemma][0][]035B}{[1][0][]source}{}{}{}}
\makeatother
\crefname{originallemma}{original lemma}{original lemmas}
\Crefname{originallemma}{Original lemma}{Original lemmas}
\makeatletter
\newlabel{properties-lemma-locally-Noetherian-quasi-separated}{{01OY}{source}{}{sourcepointer-5}{}}
\newlabel{properties-lemma-locally-Noetherian-quasi-separated@cref}{{[originallemma][0][]01OY}{[1][0][]source}{}{}{}}
\makeatother
\crefname{originallemma}{original lemma}{original lemmas}
\Crefname{originallemma}{Original lemma}{Original lemmas}
\makeatletter
\newlabel{schemes-lemma-compose-after-separated}{{01KV}{source}{}{sourcepointer-6}{}}
\newlabel{schemes-lemma-compose-after-separated@cref}{{[originallemma][0][]01KV}{[1][0][]source}{}{}{}}
\makeatother
\title{The normalization in a reduced scheme is a finite morphism over a field}
\author{The Stacks Project Authors}
\date{stacks-project:varieties.tex:local-snapshot-20261011; GNU FDL 1.2 or later}
\begin{document}
\maketitle
\noindent\textit{Editorial note.} All delivered mathematical text is the author's own, copied from the pinned Stacks Project chapter file varieties.tex; no statement or proof step was rewritten, reordered or removed.\par\medskip
\noindent\textit{Editorial note.} The proof ends with the author's own 'Some small details omitted.' after the two uses of local algebraicity; the sentence is reproduced verbatim. The hypotheses used are named explicitly, so the omission is bounded.\par\medskip
\noindent\textit{Editorial note.} This card was re-admitted after a gate audit: a text-similarity screen had flagged it against a statement of another chapter whose conclusion is that the normalization is reduced rather than finite. The two are different propositions, so the ruling recorded for this lane is to keep this card.\par\medskip
\noindent\textit{Editorial note.} The counted spans carry 6 navigation stubs and no bibliography citation.\par\medskip
\noindent\textit{Editorial note.} The preamble is the chapter's retained minimal preamble (sources/stacks-varieties/retained-preamble.tex), which since 2026-10-12 also carries the project's own reference/slogan/history definitions.\par\medskip
\noindent\textit{Allowed standard prerequisites (editorial).}\begin{itemize}
\item Locally algebraic schemes over a field are locally Noetherian and quasi-separated. Properties, Lemma properties-lemma-locally-Noetherian-quasi-separated and Morphisms, Lemma morphisms-lemma-finite-type-noetherian, invoked at proof lines 4978-4981. Supplies that Y is quasi-separated, hence that the morphism is quasi-separated.
\item Quasi-separatedness is stable under composition with a separated morphism. Schemes, Lemma schemes-lemma-compose-after-separated, invoked at proof lines 4981-4982. Supplies that the morphism f is quasi-separated.
\item Definition of the normalization of a scheme in another scheme. Morphisms, Definition morphisms-definition-normalization-X-in-Y, invoked at proof lines 4983-4984. Supplies the object X' whose finiteness is at stake.
\item Finiteness of normalization in a reduced scheme over a Nagata base. Morphisms, Lemma morphisms-lemma-nagata-normalization-finite-general and Lemma morphisms-lemma-ubiquity-nagata, invoked at proof lines 4984-4989. Supplies the finiteness conclusion.
\end{itemize}\medskip
\ifdefined\hypertarget\hypertarget{sourcepointer-1}{}\fi
\noindent\textit{Source reference.} Original definition 035H: sources/stacks-chow/tags.txt, source lines 4128--4128. Pointer outside the counted span, invoked once for context; the load-bearing argument of the counted statement is the author's proof reproduced in full in the proof span.\par\medskip
\ifdefined\hypertarget\hypertarget{sourcepointer-2}{}\fi
\noindent\textit{Source reference.} Original lemma 01T6: sources/stacks-chow/tags.txt, source lines 2363--2363. Pointer outside the counted span, invoked once for context; the load-bearing argument of the counted statement is the author's proof reproduced in full in the proof span.\par\medskip
\ifdefined\hypertarget\hypertarget{sourcepointer-3}{}\fi
\noindent\textit{Source reference.} Original lemma 0AVK: sources/stacks-chow/tags.txt, source lines 13688--13688. Pointer outside the counted span, invoked once for context; the load-bearing argument of the counted statement is the author's proof reproduced in full in the proof span.\par\medskip
\ifdefined\hypertarget\hypertarget{sourcepointer-4}{}\fi
\noindent\textit{Source reference.} Original lemma 035B: sources/stacks-chow/tags.txt, source lines 4122--4122. Pointer outside the counted span, invoked once for context; the load-bearing argument of the counted statement is the author's proof reproduced in full in the proof span.\par\medskip
\ifdefined\hypertarget\hypertarget{sourcepointer-5}{}\fi
\noindent\textit{Source reference.} Original lemma 01OY: sources/stacks-chow/tags.txt, source lines 2211--2211. Pointer outside the counted span, invoked once for context; the load-bearing argument of the counted statement is the author's proof reproduced in full in the proof span.\par\medskip
\ifdefined\hypertarget\hypertarget{sourcepointer-6}{}\fi
\noindent\textit{Source reference.} Original lemma 01KV: sources/stacks-chow/tags.txt, source lines 2064--2064. Pointer outside the counted span, invoked once for context; the load-bearing argument of the counted statement is the author's proof reproduced in full in the proof span.\par\medskip
\begin{lemma}
\label{lemma-relative-normalization-finite}
Let $k$ be a field. Let $f : Y \to X$ be a quasi-compact
morphism of locally algebraic schemes over $k$. Let $X'$
be the normalization of $X$ in $Y$. If $Y$ is reduced, then
$X' \to X$ is finite.
\end{lemma}

\begin{proof}
Since $Y$ is quasi-separated (by
Properties, Lemma \ref{properties-lemma-locally-Noetherian-quasi-separated} and
Morphisms, Lemma \ref{morphisms-lemma-finite-type-noetherian})
the morphism $f$ is quasi-separated
(Schemes, Lemma \ref{schemes-lemma-compose-after-separated}).
Hence Morphisms, Definition \ref{morphisms-definition-normalization-X-in-Y}
applies. The result follows from Morphisms, Lemma
\ref{morphisms-lemma-nagata-normalization-finite-general}.
This uses that locally algebraic schemes are locally Noetherian
(hence have locally finitely many irreducible components)
and that locally algebraic schemes are Nagata
(Morphisms, Lemma \ref{morphisms-lemma-ubiquity-nagata}).
Some small details omitted.
\end{proof}

\end{document}
